% called by main.tex
%

\section*{Abstract}
\addcontentsline{toc}{section}{Abstract}
\label{sec::abstract}

This work studies whether deep learning can exploit market microstructure signals to
anticipate the short-term price movement of binary contracts on Bitcoin traded on
Polymarket, a prediction market. Unlike the literature focused on how the event finally
resolves, here we target the local movements of the contract price over windows of a few
seconds, combining the order book with external Bitcoin market references (spot, perpetual
futures and a price oracle). The main contribution is not a single model but a
reproducible methodology: a multi-source data capture system with per-session quality
control, strictly temporal validation, and an evaluation that incorporates realistic
execution costs and latency. The central, uncomfortable finding is that a directional
classifier that is correct nearly nine out of ten times stops being profitable once a
realistic two-second entry latency is introduced: what decides the outcome is cost and
execution, not directional accuracy. Building on this, we evaluate a progression of models
(tabular baselines, latency-aware models, and sequential and convolutional architectures
over the order book) and a frozen tabular specialist. On the untouched five-day
out-of-sample block, the specialist obtains a modest positive mean at the reference cost
(+0.349 ticks per action, 754 actions and three out of five positive days), but turns
slightly negative under the full-cost stress test. A volatility filter, whose numerical
threshold came from the training period, is then explored after detecting a regime shift;
it improves the result on the same block, so it is reported as a post-hoc diagnostic rather
than as confirmatory evidence. We report every figure in net ticks together with its support
and uncertainty, and we do not claim real profitability: the deliverable is the
methodology, the evidence obtained under an auditable protocol, and the explicit
delimitation of its limits.

\newpage
\section*{Resumen}
\addcontentsline{toc}{section}{Resumen}
\label{sec::resumen}

Este trabajo estudia si el aprendizaje profundo puede aprovechar señales de
microestructura de mercado para anticipar el movimiento de corto plazo del precio de los
contratos binarios sobre Bitcoin negociados en Polymarket, un mercado de predicción. A
diferencia de la literatura centrada en cómo se resuelve finalmente el evento, aquí el
objetivo son los movimientos locales del precio del contrato en ventanas de pocos
segundos, combinando el libro de órdenes con referencias externas del mercado de Bitcoin
(contado, futuros perpetuos y un oráculo de precios). La aportación principal no es un
modelo concreto, sino una metodología reproducible: un sistema de captura de datos
multifuente con control de calidad por sesión, una validación estrictamente temporal y una
evaluación que incorpora costes de ejecución y latencia realistas. El hallazgo central, y
algo incómodo, es que un clasificador direccional que acierta casi nueve de cada diez veces
deja de ser rentable en cuanto se introduce una latencia de entrada realista de dos
segundos: lo que decide el resultado es el coste y la ejecución, no el acierto direccional.
A partir de ahí se evalúa una progresión de modelos (líneas base tabulares, modelos
sensibles a la latencia y arquitecturas secuenciales y convolucionales sobre el libro de
órdenes) y un especialista tabular congelado. En el bloque intacto de cinco días fuera de
muestra, el especialista obtiene un promedio modestamente positivo al coste de referencia
($+0{,}349$ ticks por acción, 754 acciones y tres de cinco días positivos), pero pasa a ser
ligeramente negativo bajo el coste de estrés completo. Tras detectar un cambio de régimen se
explora un filtro de volatilidad cuyo umbral numérico procede del entrenamiento; como la
decisión de incorporarlo se evalúa sobre ese mismo bloque, su mejora se presenta como
diagnóstico \emph{post-hoc}, no como evidencia confirmatoria. Todos los resultados se reportan
en ticks netos junto a su soporte e incertidumbre, y no se afirma rentabilidad real: lo que
se entrega es la metodología, la evidencia obtenida bajo un protocolo auditable y la
delimitación explícita de sus límites.
